\paragraph{Product backgrounds \(M\times K\) at the maximally supersymmetric content: factorisation of the one-loop trace and the index part of the induced Einstein coefficient (conditional; bounded no-go).}
Work in the Poisson limit with the content of~\cite{IKKT1997}, supplied as an input since the displayed action is bosonic: the ten-component vector with weight \(1\), the ghosts with weight \(-2\) and one sixteen-component Weyl spinor with weight \(-\tfrac12\), with operators \(H_0+\Sigma^{AB}\mathbb F_{AB}\), where \(H_0\) is the common scalar operator, \(\Sigma^{AB}\) are the generators of the vector or spinor representation and \(\mathbb F_{AB}\) is the action of the flux \(\{X_A,X_B\}\). Take a product background: the indices split into a block \(M\) of dimension \(d\) and a block \(K\) of dimension \(10-d\), \(\{X^a,X^i\}=0\) for \(a\) in \(M\) and \(i\) in \(K\), the fluctuations are functions on \(M\) times matrices of a fuzzy space \(K\), and the operators of the two factors commute.

\emph{Factorisation.} On one factor write \(s_0\) for the trace of \(e^{-sH_0}\), \(v\) for the trace over the vector representation of the block minus its dimension times \(s_0\) (the index part), and, for \(d=4\), \(\sigma_\pm\) and \(\tau_\pm\) for the index parts of the chiral spinors of \(\mathfrak{so}(4)\) and \(\mathfrak{so}(6)\), \(\sigma=\sigma_++\sigma_-\), \(\tau=\tau_++\tau_-\). Then, exactly,
\[
  \operatorname{str}e^{-sH}=s_0^K\bigl(v^M-2\sigma^M\bigr)+s_0^M\bigl(v^K-\tau^K\bigr)-\tfrac12\bigl(\sigma_+^M\tau_+^K+\sigma_-^M\tau_-^K\bigr),
\]
and for \(d=7\), with the spinor blocks \(8\) of \(\mathfrak{so}(7)\) and \(2\) of \(\mathfrak{so}(3)\), \(s_0^K(v^M-\sigma_8^M)+s_0^M(v^K-4\tau_2^K)-\tfrac12\sigma_8^M\tau_2^K\) (checked with random non-commuting finite-dimensional operators, relative residual below \(10^{-12}\)). The index parts of the vector and of the spinors agree at second order in the flux, so the first two terms start at fourth order in the flux of their block, and the mixed term starts at second order in each, where it equals \(-\tfrac18v_2^Mv_2^K\) with \(v_2\) the second-order index part of the vector (Fourier extraction of the coefficients in both splits). The vector representation does not enter the mixed term, since \(\Sigma^{ab}\Sigma^{ij}=0\) on it. For \(d=4\) the mixed term is \(-\tfrac14\sigma^M\tau^K-\tfrac14(\sigma_+^M-\sigma_-^M)(\tau_+^K-\tau_-^K)\); the second part is odd under the orientation of each block, starts at third order in the flux of \(K\) and is not treated, and \(\tau^K=v_2^K\) at second order. The mixed term corresponds to the term \(2F^{\mu\nu}F_{\mu\nu}F^{ij}F_{ij}\) of Section~6.2 of~\cite{Steinacker2023}. This covers, for product backgrounds, transverse matrices with flux, which the paragraph on the supertrace at an arbitrary matrix fluctuation does not treat.

\emph{Two-derivative part for four coordinates on \(M\).} Let \(M\) be the four-coordinate class of the Presentation~A paragraphs above, \(d=4\), Euclidean section. There \(s_0^M\) has the density \(\tfrac16\rho R[g_\rho]\), and the second-order index parts of the vector and of a Dirac field have the densities \(\rho^2X\) and \(\tfrac12\rho^2X\) (paragraph on fermionic content in Section~\ref{sec:status-B}). The first term has no two-derivative part, and the part of the induced action with two derivatives on \(M\) that depends on \(K\) is, apart from the orientation-odd remainder,
\[
  -\frac1{32\pi^2}\int\Bigl(\tfrac16W_R\,\rho R[g_\rho]+W_X\,\rho^2X\Bigr),\qquad
  W_R=\int_0^\infty\frac{\mathrm ds}{s^2}\bigl(v^K-\tau^K\bigr)(s),\qquad
  W_X=-\frac18\int_0^\infty\frac{\mathrm ds}{s^2}\,\tau^K(s),
\]
both finite without a cutoff and set by the scale of \(K\) when the vector operator of \(K\) has no negative mode (the spinor operator is a square); at second order in the flux of \(K\), \(\tau^K=v_2^K\). The spinor generators are traceless, so \(\tau^K\ge0\) by the Peierls--Bogoliubov inequality (checked at random data), and \(W_X<0\) for a nonzero flux of \(K\); \(W_R\) has no fixed sign. On the two modes with \(\delta\rho=0\) about constant \(\bar\Theta\) one has \(V_\mu=0\), \(\mathcal T=X/4\) and \(\tfrac32S_{\lambda\mu\nu}S^{\lambda\mu\nu}=X/2\) identically in the polarisation (exact rational arithmetic, four backgrounds), so \(\rho^2X\) acts as \(-4\rho R[g_\rho]\), of which the term \(-2\rho R[g_\rho]\) of the presentation \(r=-2\) of the paragraph on local frame rotations is one half and the axial torsion the other; the axial part of the torsion is also nonzero on the on-shell null mode with \(\delta\rho=0\) of the Lorentzian section. In the units \(c\) of the paragraph on the order \(\Lambda^2\) in Section~\ref{sec:status-B}, with \(\Lambda_{\mathrm{UV}}^2\) replaced by these integrals (\(1/G=c/2\pi\)), the tensor coefficient is \(c_T=\tfrac16W_R-4W_X\); the mixed part \(-4W_X>0\) has the sign of a scalar component. In the presentation \(r=-2\) the coefficient of \(\rho R[g_\rho]\) is \(\tfrac16W_R-2W_X\). The coefficient that~\cite{Steinacker2023} identifies with the inverse Newton constant corresponds to its part \(-2W_X\); the contribution of \(K\) alone is evaluated in Section~6.1 there as an effective potential, and its curvature term \(\tfrac16W_R\), of the same order in the scale of \(K\), is not given there.

\emph{Area coefficient.} The cone on \(M\) leaves \(K\) a spectator, so the factorisation holds on the replica cover. With the cone values of Section~\ref{sec:status-B} (Dirac index part \(-1\) per Dirac field for the antiperiodic gluing, \(0\) for the periodic one, and \(-\tfrac12\) times the coefficient of \(\rho R[g_\rho]\), with \(X\) kept as the torsion invariant, for the smooth replica), the fixed-state replica gives \(c_A=\tfrac16W_R-2W_X\) plus \(2\) per matrix mode of \(K\) at order \(\LUV^2\) (the value \(c_A=2\) at \(N_s=6\), \(n_D=2\), carried by the first term), the periodic replica \(c_A=\tfrac16W_R\) and the smooth replica \(c_A=-\tfrac1{12}W_R\). Hence \(c_A-c_T\) is \(4W_X\) for the periodic and \(-\tfrac14W_R+4W_X\) for the smooth replica, and for the fixed-state replica a cutoff-dependent term of order \(\LUV^2\) plus the finite part \(2W_X\): for \(W_X\neq0\) the periodic gap does not vanish and the smooth one vanishes only for \(W_R=16W_X\), a condition on \(K\). The vector index part, which with the weight \(\kappa_X\) carries the gap \(4\kappa_X\) of the paragraph on the normalisation, cancels against the Dirac fields in the first term and is absent from the mixed term. The finite part \(2W_X\) of the fixed-state gap is the axial-torsion part of the Dirac index density on the tensor modes; the Dirac cone value \(-1\) equals the coefficient of \(R\) in \(\tfrac12\rho^2X\) in the presentation \(r=-2\) (an equality of the two values, not derived), so, with the same index density and cone value, the finite gap would vanish for tensor modes that source neither the axial nor the vector torsion at first order.

\emph{Classes with six polarisations.} For a product the linearised Einstein rank is that of \(M\) (paragraph on leaf dimension above four), so \(M\) must carry the six polarisations. For \(D\) matrix coordinates on a four-dimensional leaf with constant \(\theta\) in the chart \(y\), the scalar operator is \(-\lambda\Delta_{\hat g}\) with \(\hat g=\lambda G\), \(\lambda=\sqrt{\lvert g\rvert}/\rho\) and \(\lvert\det\hat g\rvert=\lvert\det g\rvert\), and its density is \(\tfrac16\rho R[\hat g]\) pointwise (exact algebra in both signatures; the density checked symbolically at a random Euclidean jet); \(\rho\) is constant, so the scalar part of the two-derivative density in this class is \(\rho\) times the Einstein--Hilbert functional of \(\hat g\). At the frozen data of the paragraph on more than four matrix coordinates the linearised Einstein operator of \(\hat g\) has rank \(\min(6,D-1)\) on the first-order image (one exact-rational draw for each \(4\le D\le11\)), and the Hamiltonian gauge direction is a linearised diffeomorphism of \(\hat g\). With \(d=7\) (four tangent and three normal directions) and \(K\) in the remaining three, the index density of the mixed term is \(X_7=g^{\rho\lambda}\partial_\rho\Theta^{IJ}\partial_\lambda\Theta_{IJ}\), which on a transverse wave \(\delta X^i=\phi_i\cos(p\cdot y)\) about \(x^i=\tfrac12\mathrm{II}^i_{\mu\nu}y^\mu y^\nu\) contains the term \((p\cdot g^{-1}p)(p\cdot Kp)\lvert\phi\rvert^2\) of fourth order in \(p\) at \(y=0\) (exact check at one background), whereas the Einstein--Hilbert form of \(\hat g\) on those modes is of second order in \(p\) by derivative counting, since \(\delta\hat g\) is algebraic in \(\phi\) at first order. The two forms are not proportional, so for a nonzero flux of \(K\) no single \(G_T\) covers the transverse polarisations at the frozen-coefficient leading symbol. The conjugate-coordinate class with \(K\) in the transverse directions has the structure \(S^2_n\,\bar\times\,M^{3,1}\times K\) of Section~7 of~\cite{Steinacker2023}, with frame \(\{t^\alpha,x^\mu\}\), where \(t^\alpha\) and \(x^\mu\) are the semiclassical forms of the generators \(T^a=R^{-1}M^{a4}\) and \(X^a=rM^{a5}\) of~\cite{SperlingSteinacker2019}, that is \(T^a\) and \(y^\mu\) here up to normalisation; there the mixed term is the spinor index part of the Dirac operator of the background \(T^a\) on \(\mathbb{CP}^{1,2}\). Its two-derivative density, its values on the six level-one amplitudes and on the cone, and whether those amplitudes source the axial torsion, are not computed; in the Poisson limit its heat trace is the hypoelliptic one of the paragraph on the \(\Theta\)-built heat kernel, whereas the cited construction uses a finite fibre.

These statements are for the Poisson limit on \(M\), one loop, the content stated, product backgrounds without mixed flux, the orientation-even part of the mixed term, the replicas of Section~\ref{sec:status-B} with the hypotheses stated there, and, for the transverse modes, frozen coefficients. They do not construct \(K\) or fix the sign of \(W_R\), so no value or sign of \(G_{\mathrm{eff}}\) follows; they do not derive the replica selected by the matrix state, identify \(G\) or \(\hat g\) with the emergent metric of Chapter~\ref{ch:vierbein-emergence}, or treat Moyal corrections, non-product backgrounds or a finite fibre. Problems~B and~F remain open and PROGRAMME-CONDITIONAL.
